\documentclass[]{beamer}
\usepackage[utf8]{inputenc}
\usepackage[english]{babel}
\usepackage{hyperref}
\usepackage{ulem}\normalem
\usepackage{textcomp}

\newcommand{\fondoTitulo}{%
  \usebackgroundtemplate{\includegraphics[width=\paperwidth]{fondo}}}
\newcommand{\fondoTexto}{%
  \usebackgroundtemplate{\includegraphics[width=\paperwidth]{fondo}}}

\definecolor{celeste}{HTML}{5E91AA}
\definecolor{azul}{HTML}{163F54}

\setbeamercolor{head1}{fg=celeste}
\setbeamercolor{title}{fg=celeste}
\setbeamercolor{subtitle}{fg=celeste}
\setbeamercolor{frametitle}{fg=celeste}
\setbeamercolor{structure}{fg=azul}
\setbeamercolor{normal text}{fg=azul}

\title[jMock->Mockito]{\textbf{\LARGE{Migrando de jMock a Mockito}}}
\subtitle[]{Evoluciona o muere}
\author[aranax]{Wences Arana}
\institute[lugusac, debian-gt, guatejug]{lugusac, debian-gt, guatejug}
\date[14/11/2015]{14 de noviembre de 2015}

\AtBeginSection[]
{
  \begin{frame}
    \frametitle{Programa}
    \tableofcontents[currentsection]
  \end{frame}
}

\begin{document}

\fondoTitulo
\begin{frame}
  \titlepage
\end{frame}

\fondoTexto
\transboxout
\section{Frameworks}
\subsection{jMock}
\begin{frame}[c]
  \frametitle{jMock}
  \textbf{Features}\footnote{http://www.jmock.org/index.html}
  \pause
  \begin{itemize}
    \item Provee una extensa libreria de mocking objects
    \item Fácil y rápido
    \item Librería extensea
    \item Completamente declarativa
    \item Fácil de extender
    \item expect-run-verify
  \end{itemize}
\end{frame}

\subsection{Mockito}
\begin{frame}
  \frametitle{Mockito}
  \textbf{Features:}\footnote{https://code.google.com/p/mockito/wiki/FeaturesAndMotivations}
  \begin{itemize}
  \pause \item Desarrollo continuo (push)
  \pause \item JDK8 compatible (Defenders)
  \pause \item API limpia y simple
  \pause \item Fácil lectura
  \pause \item Bastante utilizada
  \end{itemize}
\end{frame}

\begin{frame}
  \frametitle{Mockito}
  \textbf{Features(continuación...)}
  \begin{itemize}
  \pause \item Mockea clasess además de interfaces
  \pause \item ANNOTATIONS o/ 
  \pause \item Stack trace limpio 
  \end{itemize}
\end{frame}

\begin{frame}[c]
  \frametitle{Cómo funciona?}
  \begin{itemize}
  \pause \item \alert{No es una libreria expect-run-verify}
  \pause \item \alert{NO EXISTEN LAS EXPECTATIONS}
  \pause \item Preguntar sobre interacciones despues de ejecutar
  \pause \item Verifica solo lo que necesites (I can test what I want)
  \pause \item Stubbing va antes de la ejecución, verificaciones de interacción después
  \end{itemize}
\end{frame}

\section{Diferencias entre jMock y Mockito}
\begin{frame}
\frametitle{Mockito vs jMock}
  \begin{itemize}
  \item Jmock usa un contexto para manejar los mocks
  \item Mockito lo hace con metodos importados estáticamente
  \item jMock esta diseñado para mockear interfaces, no clases
  \end{itemize}
\end{frame}

\begin{frame}{Mockito vs Jmock}
  \begin{itemize}
  \item Mockito utiliza multiples mocks para la misma clase por medio de identificadores
    \item section matchers
    \item uso de any
    \item ambos tienen problemas combinando
  \item numero de invocaciones
    \item mockito usa verify
    \item mockito no revisa si se hace una llamada no esperada
    \item jmock tiene allowing e ignoring
  \item call order
  \item jmock usa sequence
    \item mockito usa inOrder
  \end{itemize}

\end{frame}

%team

%game

%Review state verification _after the tests have run_
%Behavior verification _checking to see if certain methods where invoked(matchers, etc)_

%JMOCk: set up expectations, then execute and finally verify
%Mockito verify mock invocation after the test is executed

%Ver http://blog.trafficparrot.com/2014/05/do-you-have-problems-with-your-test.html
%https://www.linkedin.com/pulse/what-difference-between-stub-mock-virtual-service-wojciech-bulaty

%Spaceballs
%https://www.linkedin.com/pulse/what-difference-between-stub-mock-virtual-service-wojciech-bulaty

%http://www.vogella.com/tutorials/Mockito/article.html

%ohttp://jmockit.org/MockingToolkitComparisonMatrix.html

%http://blogs.justenougharchitecture.com/unit-testing-with-mocks-easymock-jmock-and-mockito/


%% \fondoTitulo
\begin{frame}[c]
   \begin{center}
     \LARGE ¡Gracias!
   \end{center}    
\end{frame}

\end{document}
